\begin{Example}
\label{ex:fourier-series-distribution}
In this and the following example, we exploit the fact that differentiation is a continuous operation. Let us consider the Fourier series
\begin{equation}\label{eq:cosine-series-distributional-sum}
\cos x+\cos2x+\cdots+\cos nx+\cdots.
\end{equation}
Clearly, this series does not converge in the ordinary sense; for example, it diverges for $x=0$. However, the series
\begin{equation}\label{eq:sawtooth-sine-series}
\sin x+\dfrac12\sin2x+\dfrac13\sin3x+\cdots
\end{equation}
converges to $(\pi-x)/2$ on $(0,2\pi)$, uniformly on every compact subinterval, and the partial sums of the series \eqref{eq:sawtooth-sine-series} are uniformly bounded on $\mathbb{R}$. (We shall not prove these claims here; instead we refer to texts on Fourier series or advanced calculus or to the discussion of Fourier series in Chapter~\ref{chap:function-spaces}.) From this, it is clear that \eqref{eq:sawtooth-sine-series} converges in the sense of distributions to the $2\pi$-periodic continuation of $(\pi-x)/2$; that is, a "sawtooth wave" with jumps at integer multiples of $2\pi$. We can even write this in terms of the Heaviside function.\[
\begin{aligned}
&\sin x+\dfrac12\sin2x+\dfrac13\sin3x+\cdots\\
&\quad=\dfrac{\pi-x}2+\pi\sum_{n=1}^\infty H(x-2\pi n)-\pi\sum_{n=0}^\infty H(-2\pi n-x)
\end{aligned}
\]We now obtain \eqref{eq:cosine-series-distributional-sum} by differentiation and the symmetry of the dirac delta:
\begin{equation}\label{eq:cosine-series-dirac-comb}
\begin{aligned}
\cos x+\cos2x+\cos3x+\cdots
&=\dfrac{\mathrm{d}}{\mathrm{d}x}\left(\sin x+\dfrac12\sin2x+\dfrac13\sin3x+\cdots\right)\\
&=-\dfrac12+\pi\sum_{n\in\mathbb{Z}}\delta(x-2\pi n)
\end{aligned}
\end{equation}
(cf. Problem~\ref{ex:piecewise-smooth-distribution-derivative}).
\end{Example}

\begin{Example}
\label{ex:delta-sequence-primitive}
To prove that a sequence of integrable functions $f_n:\mathbb{R}\to\mathbb{R}$ converges to the delta function, it suffices to show that the primitives converge to the Heaviside function. The following conditions are sufficient for this:
\begin{enumerate}\item For any $\epsilon>0$, we have
\begin{equation}\label{eq:delta-sequence-uniform-tail-integrals}
\lim_{n\to\infty}\int_{-\infty}^{-a}f_n(x)\,\mathrm{d}x=0,\quad\lim_{n\to\infty}\int_a^\infty f_n(x)\,\mathrm{d}x=0
\end{equation}
uniformly for $a\in[\epsilon,\infty)$;
\item 
\begin{equation}\label{eq:delta-sequence-unit-integral-limit}
\lim_{n\to\infty}\int_{-\infty}^\infty f_n(x)\,\mathrm{d}x=1;
\end{equation}

\item $\left|\int_{-\infty}^af_n(x)\,\mathrm{d}x\right|$ is bounded by a constant independent of $a\in\mathbb{R}$ and $n\in\mathbb{N}$.
\end{enumerate}
Examples of functions satisfying these conditions are
\begin{equation}\label{eq:cauchy-kernel-delta-sequence}
f_\epsilon(x)=\dfrac{\epsilon}{\pi(x^2+\epsilon^2)},\quad\epsilon\to0,
\end{equation}

\begin{equation}\label{eq:heat-kernel-delta-sequence}
f_t(x)=\dfrac1{2\sqrt{\pi t}}\exp\left(-\dfrac{x^2}{4t}\right),\quad t\to0^+,
\end{equation}

\begin{equation}\label{eq:sinc-kernel-delta-sequence}
f_n(x)=\dfrac{\sin nx}{\pi x},\quad n\to\infty.
\end{equation}

\end{Example}
